\subsection{Enfoque de Evaluación de Resultados}\label{subsec:enfoque-resultados}

La naturaleza del proyecto —una intervención de ingeniería sobre una plataforma en producción con base de usuarios activa— determina el tipo de evidencia con que se valoran sus resultados. En lugar de contrastes estadísticos o estudios con usuarios, la evaluación se apoya en tres clases de evidencia complementarias: (i) la verificación de artefactos, que constata la existencia y el funcionamiento de los componentes construidos; (ii) la confrontación entre el resultado esperado declarado en el alcance y el resultado efectivamente obtenido para cada objetivo específico; y (iii) la demostración mediante artefactos concretos, tales como inspecciones del tráfico de red, trazas de telemetría correlacionadas, paneles de monitoreo y registros del proceso de integración continua.

Cada uno de los apartados siguientes aborda un objetivo específico, recupera el resultado esperado y expone la evidencia del resultado alcanzado. El capítulo cierra con una tabla de síntesis que consolida la confrontación objetivo por objetivo. La validación basada en pruebas automatizadas —tanto las pruebas unitarias del \textit{backend} desacoplado como la suite de reglas de seguridad de Firestore— se documenta de forma específica en el capítulo dedicado a las pruebas del sistema; en este capítulo se la referencia como evidencia de apoyo sin reproducir su detalle.
